\documentclass[10pt,a4paper]{amsart}
\usepackage{amsmath,amssymb,amsthm,graphicx}
\usepackage{verbatim}
\newenvironment{reference}{\comment}{\endcomment}
%\newenvironment{reference}{}{}
\newenvironment{slogan}{\comment}{\endcomment}
\newenvironment{history}{\comment}{\endcomment}

% For commutative diagrams we use Xy-pic
\usepackage[all]{xy}

% We use 2cell for 2-commutative diagrams.
\xyoption{2cell}
% Theorem environments.
%
\theoremstyle{plain}
\newtheorem{theorem}[subsection]{Theorem}
\newtheorem{proposition}[subsection]{Proposition}
\newtheorem{lemma}[subsection]{Lemma}

\theoremstyle{definition}
\newtheorem{definition}[subsection]{Definition}
\newtheorem{example}[subsection]{Example}
\newtheorem{exercise}[subsection]{Exercise}
\newtheorem{situation}[subsection]{Situation}

\theoremstyle{remark}
\newtheorem{remark}[subsection]{Remark}
\newtheorem{remarks}[subsection]{Remarks}

\numberwithin{equation}{subsection}

% Macros
%
\def\lim{\mathop{\mathrm{lim}}\nolimits}
\def\colim{\mathop{\mathrm{colim}}\nolimits}
\def\Spec{\mathop{\mathrm{Spec}}}
\def\Hom{\mathop{\mathrm{Hom}}\nolimits}
\def\Ext{\mathop{\mathrm{Ext}}\nolimits}
\def\SheafHom{\mathop{\mathcal{H}\!\mathit{om}}\nolimits}
\def\SheafExt{\mathop{\mathcal{E}\!\mathit{xt}}\nolimits}
\def\Sch{\mathit{Sch}}
\def\Mor{\mathop{\mathrm{Mor}}\nolimits}
\def\Ob{\mathop{\mathrm{Ob}}\nolimits}
\def\Sh{\mathop{\mathit{Sh}}\nolimits}
\def\NL{\mathop{N\!L}\nolimits}
\def\CH{\mathop{\mathrm{CH}}\nolimits}
\def\proetale{{pro\text{-}\acute{e}tale}}
\def\etale{{\acute{e}tale}}
\def\QCoh{\mathit{QCoh}}
\def\Ker{\mathop{\mathrm{Ker}}}
\def\Im{\mathop{\mathrm{Im}}}
\def\Coker{\mathop{\mathrm{Coker}}}
\def\Coim{\mathop{\mathrm{Coim}}}

% Boxtimes
%
\DeclareMathSymbol{\boxtimes}{\mathbin}{AMSa}{"02}

%
% Macros for moduli stacks/spaces
%
\def\QCohstack{\mathcal{QC}\!\mathit{oh}}
\def\Cohstack{\mathcal{C}\!\mathit{oh}}
\def\Spacesstack{\mathcal{S}\!\mathit{paces}}
\def\Quotfunctor{\mathrm{Quot}}
\def\Hilbfunctor{\mathrm{Hilb}}
\def\Curvesstack{\mathcal{C}\!\mathit{urves}}
\def\Polarizedstack{\mathcal{P}\!\mathit{olarized}}
\def\Complexesstack{\mathcal{C}\!\mathit{omplexes}}
% \Pic is the operator that assigns to X its picard group, usage \Pic(X)
% \Picardstack_{X/B} denotes the Picard stack of X over B
% \Picardfunctor_{X/B} denotes the Picard functor of X over B
\def\Pic{\mathop{\mathrm{Pic}}\nolimits}
\def\Picardstack{\mathcal{P}\!\mathit{ic}}
\def\Picardfunctor{\mathrm{Pic}}
\def\Deformationcategory{\mathcal{D}\!\mathit{ef}}
\usepackage{cleveref}
\crefname{originalreference}{original reference}{original references}
\Crefname{originalreference}{Original reference}{Original references}
\makeatletter
\newlabel{situation-setup-base-change}{{0FVG}{source}{}{sourcepointer-1}{}}
\newlabel{situation-setup-base-change@cref}{{[originalreference][0][]0FVG}{[1][0][]source}{}{}{}}
\makeatother
\crefname{originallemma}{original lemma}{original lemmas}
\Crefname{originallemma}{Original lemma}{Original lemmas}
\makeatletter
\newlabel{lemma-rational-equivalence-family}{{02S4}{source}{}{sourcepointer-2}{}}
\newlabel{lemma-rational-equivalence-family@cref}{{[originallemma][0][]02S4}{[1][0][]source}{}{}{}}
\makeatother
\crefname{originallemma}{original lemma}{original lemmas}
\Crefname{originallemma}{Original lemma}{Original lemmas}
\makeatletter
\newlabel{lemma-closed-subscheme-cross-p1}{{02S5}{source}{}{sourcepointer-3}{}}
\newlabel{lemma-closed-subscheme-cross-p1@cref}{{[originallemma][0][]02S5}{[1][0][]source}{}{}{}}
\makeatother
\crefname{originallemma}{original lemma}{original lemmas}
\Crefname{originallemma}{Original lemma}{Original lemmas}
\makeatletter
\newlabel{lemma-pullback-coherent-base-change}{{0FVI}{source}{}{sourcepointer-4}{}}
\newlabel{lemma-pullback-coherent-base-change@cref}{{[originallemma][0][]0FVI}{[1][0][]source}{}{}{}}
\makeatother
\crefname{originallemma}{original lemma}{original lemmas}
\Crefname{originallemma}{Original lemma}{Original lemmas}
\makeatletter
\newlabel{lemma-dimension-base-change}{{0FVH}{source}{}{sourcepointer-5}{}}
\newlabel{lemma-dimension-base-change@cref}{{[originallemma][0][]0FVH}{[1][0][]source}{}{}{}}
\makeatother
\title{The flat pullback map on cycles descends to a map on Chow groups with the expected degree shift}
\author{The Stacks Project Authors}
\date{stacks-project:chow.tex:local-snapshot-20261012; GNU FDL 1.2 or later}
\begin{document}
\maketitle
\noindent\textit{Editorial note.} All delivered mathematical text is the author's own, copied from the pinned Stacks Project chapter file chow.tex; no statement or proof step was rewritten, reordered or removed.\par\medskip
\noindent\textit{Editorial note.} Preamble provenance: the pinned chapter file opens with an input of the shared preamble and carries no preamble text of its own, so the project's single Stacks-wide preamble staged at sources/stacks-chow/preamble.tex is excerpted verbatim as two spans, lines 14-18 (Xy-pic setup) and lines 154-217 (the native theorem declarations and mathematical macros). Nothing was added, removed or redefined; the wrapper class and the amsmath/amssymb/amsthm/graphicx base packages are supplied by the assembler.\par\medskip
\noindent\textit{Editorial note.} The counted statement and proof contain no cross reference, no citation and no included file, so the delivered body is self contained.\par\medskip
\noindent\textit{Editorial note.} Canonical location of the counted proposition: Stacks tag 0FVJ (sources/stacks-chow/pinned-original-source.tex line 15255).\par\medskip
\noindent\textit{Allowed standard prerequisites (editorial).}\begin{itemize}
\item Rational equivalence exhibited by families over P\textasciicircum{}1. Same chapter, Lemma labelled lemma-rational-equivalence-family, cited at proof lines 15268-15270. Reduces the lemma to the two assertions for an integral W inside X x P\textasciicircum{}1.
\item The dimension shift under base change. Same chapter, Lemma labelled lemma-dimension-base-change, cited at proof lines 15292-15294. Gives that components of W' have delta'-dimension k + c.
\item The cycle of the special fibre and its base change. Same chapter, Lemma labelled lemma-pullback-coherent-base-change, cited at proof lines 15288-15290, and Lemma labelled lemma-closed-subscheme-cross-p1, cited at proof lines 15282-15286. Identifies the base change of the special fibre cycle and supplies the hypotheses to be checked for W'.
\item Nonzerodivisors and associated points under flat base change. The local computation at proof lines 15294-15299 using flatness of O\_\{W,w\} -> O\_\{W',w'\}. Shows W' has no associated points on W'\_0.
\end{itemize}\medskip
\ifdefined\hypertarget\hypertarget{sourcepointer-1}{}\fi
\noindent\textit{Source reference.} Original reference 0FVG: sources/stacks-chow/pinned-original-source.tex, source lines 15123--15123. Same-chapter pointer into the author environment situation at source line 15123; invoked at proof line None of the counted proof. It is retained only as a navigation pointer; every load-bearing step is inside the delivered span.\par\medskip
\ifdefined\hypertarget\hypertarget{sourcepointer-2}{}\fi
\noindent\textit{Source reference.} Original lemma 02S4: sources/stacks-chow/pinned-original-source.tex, source lines 3335--3335. Same-chapter pointer into the author environment lemma at source line 3335; invoked at proof line 15268 of the counted proof. It is retained only as a navigation pointer; every load-bearing step is inside the delivered span.\par\medskip
\ifdefined\hypertarget\hypertarget{sourcepointer-3}{}\fi
\noindent\textit{Source reference.} Original lemma 02S5: sources/stacks-chow/pinned-original-source.tex, source lines 3404--3404. Same-chapter pointer into the author environment lemma at source line 3404; invoked at proof line 15281 of the counted proof. It is retained only as a navigation pointer; every load-bearing step is inside the delivered span.\par\medskip
\ifdefined\hypertarget\hypertarget{sourcepointer-4}{}\fi
\noindent\textit{Source reference.} Original lemma 0FVI: sources/stacks-chow/pinned-original-source.tex, source lines 15193--15193. Same-chapter pointer into the author environment lemma at source line 15193; invoked at proof line 15287 of the counted proof. It is retained only as a navigation pointer; every load-bearing step is inside the delivered span.\par\medskip
\ifdefined\hypertarget\hypertarget{sourcepointer-5}{}\fi
\noindent\textit{Source reference.} Original lemma 0FVH: sources/stacks-chow/pinned-original-source.tex, source lines 15139--15139. Same-chapter pointer into the author environment lemma at source line 15139; invoked at proof line 15290 of the counted proof. It is retained only as a navigation pointer; every load-bearing step is inside the delivered span.\par\medskip
\begin{lemma}
\label{lemma-pullback-base-change}
In Situation \ref{situation-setup-base-change} let $X \to S$ be locally
of finite type and let $X' \to S'$ be the base change by $S' \to S$.
The map $g^* : Z_k(X) \to Z_{k + c}(X')$ above factors through rational
equivalence to give a map
$$
g^* : \CH_k(X) \longrightarrow \CH_{k + c}(X')
$$
of chow groups.
\end{lemma}

\begin{proof}
Suppose that $\alpha \in Z_k(X)$ is a $k$-cycle which is rationally equivalent
to zero. By Lemma \ref{lemma-rational-equivalence-family}
there exists a locally finite family of integral closed subschemes
$W_i \subset X \times \mathbf{P}^1$ of $\delta$-dimension $k$
not contained in the divisors
$(X \times \mathbf{P}^1)_0$ or $(X \times \mathbf{P}^1)_\infty$
of $X \times \mathbf{P}^1$ such that
$\alpha = \sum ([(W_i)_0]_k - [(W_i)_\infty]_k)$.
Thus it suffices to prove for $W \subset X \times \mathbf{P}^1$
integral closed of $\delta$-dimension $k$ not contained in the divisors
$(X \times \mathbf{P}^1)_0$ or $(X \times \mathbf{P}^1)_\infty$
of $X \times \mathbf{P}^1$ we have
\begin{enumerate}
\item the base change $W' \subset X' \times \mathbf{P}^1$ satisfies the
assumptions of Lemma \ref{lemma-closed-subscheme-cross-p1} with
$k$ replaced by $k + c$, and
\item $g^*[W_0]_k = [(W')_0]_{k + c}$ and
$g^*[W_\infty]_k = [(W')_\infty]_{k + c}$.
\end{enumerate}
Part (2) follows immediately from
Lemma \ref{lemma-pullback-coherent-base-change} and the fact that
$(W')_0$ is the base change of $W_0$ (by associativity of fibre products).
For part (1), first the statement on dimensions follows
from Lemma \ref{lemma-dimension-base-change}.
Then let $w' \in (W')_0$ with image $w \in W_0$
and $z \in \mathbf{P}^1_S$. Denote $t \in \mathcal{O}_{\mathbf{P}^1_S, z}$
the usual equation for $0 : S \to \mathbf{P}^1_S$.
Since $\mathcal{O}_{W, w} \to \mathcal{O}_{W', w'}$ is flat
and since $t$ is a nonzerodivisor on $\mathcal{O}_{W, w}$
(as $W$ is integral and $W \not = W_0$) we see that also
$t$ is a nonzerodivisor in $\mathcal{O}_{W', w'}$. Hence
$W'$ has no associated points lying on $W'_0$.
\end{proof}

\end{document}
